\documentclass[aspectratio=169]{beamer}

% ==== Theme & basic look ====
\usetheme{Madrid}
\usecolortheme{seahorse}
\setbeamertemplate{navigation symbols}{}
\setbeamertemplate{footline}[frame number]

% Packages for mathematics
\usepackage{amsmath}
\usepackage{amsfonts}
\usepackage{amssymb}
\usepackage{mathtools}
\usepackage{booktabs}
\usepackage{array}
\usepackage{tikz}
\usepackage{CJKutf8}

% Proof environment for beamer, consistent with lecture1
\providecommand{\proofname}{Proof}
\newenvironment<>{proofs}[1][\proofname]{%
    \par
    \def\insertproofname{#1.}%
    \usebeamertemplate{proof begin}#2}
  {\usebeamertemplate{proof end}}
\makeatother

% ---------- Math helpers ----------
\DeclareMathOperator{\Cat}{Cat}
\DeclareMathOperator{\KL}{D_{KL}}
\DeclareMathOperator{\DIS}{D_{IS}}
\DeclareMathOperator{\softmax}{softmax}
\DeclareMathOperator*{\argmax}{arg\,max}

\newcommand{\Prob}{\mathrm{P}}
\newcommand{\E}{\mathbb{E}}
\newcommand{\dd}{\,\mathrm{d}}
\newcommand{\cV}{\mathcal{V}}
\newcommand{\one}{\mathbf{1}}
\newcommand{\mvec}{\mathbf{m}}
\newcommand{\uvec}{\mathbf{u}}
\newcommand{\xvec}{\mathbf{x}}
\newcommand{\Rhat}{\widehat R}
\newcommand{\DeltaV}{\Delta^{|\cV|-1}}
\newcommand{\const}{\mathrm{const}}

\newcolumntype{L}[1]{>{\raggedright\arraybackslash}p{#1}}
\newcolumntype{C}[1]{>{\centering\arraybackslash}p{#1}}

\title[Discrete Diffusion]{Lecture 7: Formal Discrete Diffusion}
\subtitle{From CTMC reverse rates to SEDD and GIDD}
\author{Pipi Hu}
\institute[BIMSA]{Beijing Institute of Mathematical Sciences and Applications}
\date{\today}

\pgfdeclareimage[width=2cm]{logo}{../figures/logo.png}
\logo{\pgfuseimage{logo}}

\setbeamertemplate{navigation symbols}{}

\AtBeginSection[]
{
  \begin{frame}<beamer>{Outline}
    \tableofcontents[currentsection]
  \end{frame}
}

\AtBeginSubsection[]
{
  \begin{frame}<beamer>{Outline}
    \tableofcontents[currentsection,currentsubsection]
  \end{frame}
}

\begin{document}

\begin{frame}
  \titlepage
\end{frame}

\begin{frame}{Roadmap}
\small
\begin{block}{Logical dependency}
The lecture has one backbone:
\[
\text{forward CTMC}
\Longrightarrow
\text{true reverse rate}
\Longrightarrow
\begin{cases}
\text{score entropy for SEDD},\\
\text{CT-ELBO likelihood view}.
\end{cases}
\]
\end{block}

\begin{columns}[T,totalwidth=\textwidth]
\column{0.50\textwidth}
\begin{alertblock}{Part I: Shared mathematics}
\begin{enumerate}
  \item column-vector CTMC notation
  \item small-time reverse posterior
  \item continuous-time ELBO
\end{enumerate}
\end{alertblock}

\column{0.46\textwidth}
\begin{block}{Part II: Model families}
\begin{enumerate}
  \item SEDD: score entropy for probability ratios
  \item GIDD: interpolating marginals and posterior denoising
  \item training, sampling, and comparison
\end{enumerate}
\end{block}
\end{columns}
\end{frame}

\begin{frame}{Outline}
  \tableofcontents
\end{frame}

% ======================================================================
\section{CTMC Backbone}
% ======================================================================

\begin{frame}{Finite-state CTMC: column-vector convention}
\small
Let one token take values in a finite set
\[
  \cV=\{1,\ldots,N\}.
\]
We write probability distributions as column vectors. Matrix entries use conditional-probability order:
\[
R_t(y,x)\quad\text{is the transition rate from source }x\text{ to destination }y.
\]
Thus, for $y\neq x$,
\[
R_t(y,x)\ge 0,
\qquad
R_t(x,x)=-\sum_{y\neq x}R_t(y,x).
\]

\begin{block}{Small-time transition}
\[
\Prob(X_{t+\Delta}=y\mid X_t=x)
=\delta_{y,x}+\Delta R_t(y,x)+o(\Delta).
\]
\end{block}

\begin{alertblock}{Forward equation}
\[
\dot p_t(y)=\sum_{x\in\cV}R_t(y,x)p_t(x),
\qquad
\dot p_t=R_tp_t.
\]
\end{alertblock}
\end{frame}

\begin{frame}{Proof of the forward equation: expansion}
\scriptsize
\begin{proof}
Use total probability over the previous state at time $t$:
\[
p_{t+\Delta}(y)
=\sum_{x\in\cV}\Prob(X_{t+\Delta}=y\mid X_t=x)p_t(x).
\]
Substitute the small-time transition formula:
\[
p_{t+\Delta}(y)
=\sum_{x\in\cV}
\left(\delta_{y,x}+\Delta R_t(y,x)+o(\Delta)\right)p_t(x).
\]
Since the state space is finite, the $o(\Delta)$ terms can be summed:
\[
p_{t+\Delta}(y)
=p_t(y)+\Delta\sum_{x\in\cV}R_t(y,x)p_t(x)+o(\Delta).
\]
\end{proof}
\end{frame}

\begin{frame}{Proof of the forward equation: limit}
\scriptsize
\begin{proofs}[\proofname\ (Cont)]
Subtract $p_t(y)$, divide by $\Delta$, and let $\Delta\to0$:
\[
\frac{p_{t+\Delta}(y)-p_t(y)}{\Delta}
=\sum_{x\in\cV}R_t(y,x)p_t(x)+\frac{o(\Delta)}{\Delta}.
\]
Since $o(\Delta)/\Delta\to0$,
\[
\dot p_t(y)=
\lim_{\Delta\to0}\frac{p_{t+\Delta}(y)-p_t(y)}{\Delta}
=\sum_{x\in\cV}R_t(y,x)p_t(x).
\]
Collecting all $y\in\cV$ gives the column-vector form
\[
\dot p_t=R_tp_t.
\]
\end{proofs}
\end{frame}

\begin{frame}{When the generator has a fixed shape}
\scriptsize
Many discrete diffusion models choose
\[
R_t=\dot\sigma(t)G,
\qquad
\sigma(t)=\int_0^t\dot\sigma(s)\,\dd s.
\]
If $G$ is time-independent, then
\[
p_t=\exp(\sigma(t)G)p_0.
\]
The transition matrix over $[0,t]$ is $\exp(\sigma(t)G)$.

\begin{block}{Interpretation}
$G$ chooses the allowed token changes. The scalar schedule $\dot\sigma(t)$ is an isotropic speed: it scales all rates in $G$ by the same amount.
\end{block}

\begin{alertblock}{Convention}
The first matrix index is always the destination state and the second index is the source state:
\[
R_t(y,x)\quad\text{is the rate from }x\text{ to }y.
\]
This matches the order of the conditional probability $\Prob(X_t=y\mid X_s=x)$ and keeps the forward equation as $\dot p_t=R_tp_t$.
\end{alertblock}
\end{frame}

\begin{frame}{Sequence state space}
\small
For a length-$L$ sequence,
\[
X=(X^{(1)},\ldots,X^{(L)})\in\cV^L.
\]
The forward corruption is usually token-wise independent:
\[
q_t(Z_t\mid X)=\prod_{i=1}^Lq_t\!\left(Z_t^{(i)}\mid X^{(i)}\right).
\]

\begin{block}{Why this still models context}
The corruption process factorizes, but the denoising network sees the whole noisy sequence $Z_t$. It can predict each position using all other positions as context.
\end{block}

\begin{alertblock}{Locality of the CTMC}
With token-wise generators, a single CTMC jump changes one token. Hence the reverse rate only needs probability ratios between sequences at Hamming distance one.
\end{alertblock}
\end{frame}

\begin{frame}{Reverse posterior: keep the time labels}
\small
Fix a clean datum $x_0$ and write $q_t(z)=q_t(z\mid x_0)$.
Let
\[
s=t-\Delta,\qquad X_s=y,\qquad X_t=z.
\]
Bayes' rule gives the true one-step reverse posterior:
\[
q_{s|t}(y\mid z,x_0)
=\frac{\Prob(X_t=z\mid X_s=y)q_s(y\mid x_0)}{q_t(z\mid x_0)}.
\]

\begin{alertblock}{Why the time labels matter}
$y$ lives at the earlier time $s=t-\Delta$, while $z$ lives at the later time $t$. The term $q_s(y)$ must be Taylor expanded around time $t$:
\[
q_s(y)=q_{t-\Delta}(y)=q_t(y)-\Delta\dot q_t(y)+o(\Delta).
\]
\end{alertblock}
\end{frame}

\begin{frame}{Reverse posterior: off-diagonal case}
\scriptsize
\begin{proof}
Assume the earlier state differs from the current state: $y\neq z$. Then
\[
\Prob(X_t=z\mid X_{t-\Delta}=y)=\Delta R_t(z,y)+o(\Delta).
\]
Using the Taylor expansion of the earlier marginal,
\[
q_{t-\Delta}(y)=q_t(y)-\Delta\dot q_t(y)+o(\Delta).
\]
Therefore
\[
q_{t-\Delta|t}(y\mid z,x_0)
=
\frac{(\Delta R_t(z,y)+o(\Delta))(q_t(y)-\Delta\dot q_t(y)+o(\Delta))}
{q_t(z)}.
\]
The derivative term is multiplied by $\Delta$, so it contributes only $O(\Delta^2)$:
\[
q_{t-\Delta|t}(y\mid z,x_0)
=\Delta R_t(z,y)\frac{q_t(y)}{q_t(z)}+o(\Delta).
\]
\end{proof}
\end{frame}

\begin{frame}{Reverse posterior: diagonal case}
\scriptsize
\begin{proofs}[\proofname\ (Cont)]
Now the earlier and current state are the same token value: $y=z$. Then
\[
\Prob(X_t=z\mid X_{t-\Delta}=z)=1+\Delta R_t(z,z)+o(\Delta),
\]
and
\[
q_{t-\Delta}(z)=q_t(z)-\Delta\dot q_t(z)+o(\Delta).
\]
Substitute into Bayes' rule:
\[
q_{t-\Delta|t}(z\mid z,x_0)
=
\frac{(1+\Delta R_t(z,z)+o(\Delta))(q_t(z)-\Delta\dot q_t(z)+o(\Delta))}
{q_t(z)}.
\]
Expanding to first order,
\[
q_{t-\Delta|t}(z\mid z,x_0)
=1+\Delta\left(R_t(z,z)-\frac{\dot q_t(z)}{q_t(z)}\right)+o(\Delta).
\]
\end{proofs}
\end{frame}

\begin{frame}{True reverse rates}
\small
The off-diagonal expansion defines the true reverse jump rate from the current state $z$ at time $t$ to the earlier state $y$ at time $t-\Delta$:
\[
a_t(y,z):=R_t(z,y)\frac{q_t(y\mid x_0)}{q_t(z\mid x_0)},
\qquad y\neq z.
\]
The order is destination first: $a_t(y,z)$ is the reverse-time rate from source $z$ to destination $y$.
The diagonal coefficient is
\[
a_t(z,z):=R_t(z,z)-\frac{\dot q_t(z)}{q_t(z)}.
\]

\begin{block}{Use the forward equation}
\[
\dot q_t(z)=R_t(z,z)q_t(z)+\sum_{y\neq z}R_t(z,y)q_t(y).
\]
Thus
\[
a_t(z,z)
=-\sum_{y\neq z}R_t(z,y)\frac{q_t(y)}{q_t(z)}
=-\sum_{y\neq z}a_t(y,z).
\]
\end{block}
\end{frame}

\begin{frame}{A unified isotropic token generator}
\small
Let $\pi$ be a chosen prior distribution on $\cV$. Define the matrix
\[
\Pi_\pi(y,x)=\pi(y),
\]
so every column of $\Pi_\pi$ equals the column vector $\pi$. The isotropic redraw generator is
\[
G_\pi=\Pi_\pi-I.
\]

\begin{block}{Two important choices}
\begin{itemize}
  \item Uniform corruption: $\pi(b)=1/N$.
  \item Absorbing MASK corruption: $\pi=\delta_m$.
\end{itemize}
\end{block}

In both cases, $R_t=\dot\sigma(t)G_\pi$ and
\[
\text{the transition matrix over }[0,t]\text{ is }\exp(\sigma(t)G_\pi).
\]
\end{frame}

\begin{frame}{Why the matrix exponential simplifies}
\scriptsize
Start from a clean token $x_0$, represented by the column one-hot vector $e_{x_0}$.
Decompose it into stationary and transient parts:
\[
e_{x_0}=\pi+(e_{x_0}-\pi).
\]
For any column vector $v$,
\[
\Pi_\pi v=\pi(\one^\top v).
\]
Hence $\Pi_\pi\pi=\pi$ and
\[
G_\pi\pi=(\Pi_\pi-I)\pi=0.
\]
Also, because $\one^\top(e_{x_0}-\pi)=0$,
\[
\Pi_\pi(e_{x_0}-\pi)=0,
\qquad
G_\pi(e_{x_0}-\pi)=-(e_{x_0}-\pi).
\]
Thus $e_{x_0}-\pi$ is a right eigenvector of $G_\pi$ with eigenvalue $-1$.
\end{frame}

\begin{frame}{Using the Taylor series of the matrix exponential}
\scriptsize
By the Taylor expansion,
\[
\exp(\sigma G_\pi)
=I+\sigma G_\pi+\frac{\sigma^2}{2!}G_\pi^2+\cdots.
\]
Apply this to the two components:
\[
\exp(\sigma G_\pi)\pi=\pi,
\]
because $G_\pi^k\pi=0$ for $k\ge1$. For the transient component,
\[
\exp(\sigma G_\pi)(e_{x_0}-\pi)
=\sum_{k=0}^\infty\frac{\sigma^k}{k!}G_\pi^k(e_{x_0}-\pi)
=\sum_{k=0}^\infty\frac{\sigma^k}{k!}(-1)^k(e_{x_0}-\pi).
\]
Therefore
\[
\exp(\sigma G_\pi)(e_{x_0}-\pi)=e^{-\sigma}(e_{x_0}-\pi).
\]
\end{frame}

\begin{frame}{Closed-form token corruption}
\small
Combining the two components gives
\[
\exp(\sigma(t)G_\pi)e_{x_0}
=\pi+e^{-\sigma(t)}(e_{x_0}-\pi).
\]
Equivalently,
\[
p_{t|0}(\cdot\mid x_0)
=e^{-\sigma(t)}e_{x_0}
+(1-e^{-\sigma(t)})\pi.
\]

\begin{block}{Interpretation}
At noise level $\sigma(t)$, the token keeps its clean identity with weight $e^{-\sigma(t)}$ and is redrawn from the prior $\pi$ with weight $1-e^{-\sigma(t)}$.
\end{block}

This single formula covers both the uniform prior and the absorbing MASK prior.
\end{frame}

% ======================================================================
\section{Continuous-time ELBO}
% ======================================================================

\begin{frame}{Forward path: Markov before conditioning}
\small
Discretize $0=t_0<t_1<\cdots<t_K=1$ and write $z_i=z_{t_i}$.
The forward noising transition itself does not depend on knowing the clean datum:
\[
q_{i|i-1}(z_i\mid z_{i-1}).
\]
For the likelihood bound of a particular datum $x$, however, the proposal is the clean-conditioned forward path:
\[
q(z_{0:K}\mid x)
=q_0(z_0\mid x)\prod_{i=1}^K q_{i|i-1}(z_i\mid z_{i-1}).
\]

\begin{alertblock}{The subtle point}
The forward chain is Markov regardless of $x$. The reverse factors used in the ELBO are posterior factors under $q(z_{0:K}\mid x)$, so they depend on $x$.
\end{alertblock}
\end{frame}

\begin{frame}{Reverse posterior only after fixing $x$}
\scriptsize
For a fixed clean datum $x$, define the one-step reverse posterior by Bayes' rule:
\[
q_{i-1|i}(z_{i-1}\mid z_i,x)
=
\frac{q_{i|i-1}(z_i\mid z_{i-1})q_{i-1}(z_{i-1}\mid x)}
{q_i(z_i\mid x)}.
\]

\begin{block}{Conditioned reverse factorization}
\[
q(z_{0:K}\mid x)
=q_K(z_K\mid x)\prod_{i=1}^Kq_{i-1|i}(z_{i-1}\mid z_i,x).
\]
\end{block}

\begin{alertblock}{Why this is the right object}
The ELBO for $\log p_\theta(x)$ uses $q(z_{0:K}\mid x)$ as the proposal. Therefore the exact reverse factors are $x$-conditioned denoising posteriors, not unconditional reverse transitions.
\end{alertblock}
\end{frame}

\begin{frame}{Proof of the conditioned reverse factorization}
\scriptsize
\begin{proof}
Substitute Bayes' rule into the proposed reverse product:
\[
q_K(z_K\mid x)\prod_{i=1}^K q_{i-1|i}(z_{i-1}\mid z_i,x)
\]
\[
=q_K(z_K\mid x)\prod_{i=1}^K
\frac{q_{i|i-1}(z_i\mid z_{i-1})q_{i-1}(z_{i-1}\mid x)}
{q_i(z_i\mid x)}.
\]
The marginal factors telescope:
\[
q_K\cdot\frac{q_0q_1\cdots q_{K-1}}{q_1q_2\cdots q_K}
=q_0,
\]
where every $q_i$ is shorthand for $q_i(z_i\mid x)$.
Thus
\[
q_K(z_K\mid x)\prod_{i=1}^Kq_{i-1|i}(z_{i-1}\mid z_i,x)
=q_0(z_0\mid x)\prod_{i=1}^Kq_{i|i-1}(z_i\mid z_{i-1})
=q(z_{0:K}\mid x).
\]
\end{proof}
\end{frame}

\begin{frame}{What is structural here?}
\small
\begin{block}{A conditional-path statement}
Once $x$ is fixed, the conditional path law $q(z_{0:K}\mid x)$ has two equivalent factorizations:
\[
q_0(z_0\mid x)\prod_i q_{i|i-1}(z_i\mid z_{i-1})
\quad\text{and}\quad
q_K(z_K\mid x)\prod_i q_{i-1|i}(z_{i-1}\mid z_i,x).
\]
\end{block}

\begin{alertblock}{Modeling implication}
During training, $x$ is known, so the exact posterior $q_{i-1|i}(\cdot\mid z_i,x)$ can be computed or simplified. During generation, $x$ is unknown, so $p^\theta_{i-1|i}(\cdot\mid z_i)$ learns to replace this family of $x$-conditioned posteriors.
\end{alertblock}
\end{frame}

\begin{frame}{Start from a discrete-time ELBO}
\small
Use the exact reverse factorization as the variational proposal:
\[
q(z_{0:K}\mid x)=q_K(z_K\mid x)\prod_{i=1}^K
q_{i-1|i}(z_{i-1}\mid z_i,x).
\]
Let the model approximate these reverse posteriors by
$p^\theta_{i-1|i}(z_{i-1}\mid z_i)$.

\begin{block}{Discrete ELBO}
\[
\log p_\theta(x)\ge C_{\log}
-\sum_{i=1}^K
\E_{z_i\sim q_i(\cdot\mid x)}
\KL\!\left(
q_{i-1|i}(\cdot\mid z_i,x)
\|p^\theta_{i-1|i}(\cdot\mid z_i)
\right).
\]
\end{block}

\[
C_{\log}=\E_{q_0(z_0\mid x)}[\log p_\theta(x\mid z_0)]
-\KL(q_1(\cdot\mid x)\|p_1).
\]
\end{frame}

\begin{frame}{Discrete ELBO proof: Jensen step}
\scriptsize
\begin{proof}
Define the model path distribution
\[
p_\theta(x,z_{0:K})
=p_1(z_K)\prod_{i=1}^K
p^\theta_{i-1|i}(z_{i-1}\mid z_i)\,p_\theta(x\mid z_0).
\]
Insert the forward path distribution as a proposal:
\[
\log p_\theta(x)
=\log\sum_{z_{0:K}}q(z_{0:K}\mid x)
\frac{p_\theta(x,z_{0:K})}{q(z_{0:K}\mid x)}.
\]
By Jensen's inequality,
\[
\log p_\theta(x)
\ge
\E_{q(z_{0:K}\mid x)}
\left[
\log\frac{p_\theta(x,z_{0:K})}{q(z_{0:K}\mid x)}
\right].
\]
\end{proof}
\end{frame}

\begin{frame}{Discrete ELBO proof: collecting KL terms}
\scriptsize
\begin{proofs}[\proofname\ (Cont)]
Substitute the two path factorizations:
\[
\log\frac{p_\theta(x,z_{0:K})}{q(z_{0:K}\mid x)}
=\log p_\theta(x\mid z_0)
+\log\frac{p_1(z_K)}{q_1(z_K\mid x)}
+\sum_{i=1}^K
\log
\frac{p^\theta_{i-1|i}(z_{i-1}\mid z_i)}
{q_{i-1|i}(z_{i-1}\mid z_i,x)}.
\]
Taking expectation under $q(z_{0:K}\mid x)$ gives
\[
\E[\log p_\theta(x\mid z_0)]
-\KL(q_1(\cdot\mid x)\|p_1)
\]
\[
-\sum_{i=1}^K
\E_{z_i\sim q_i(\cdot\mid x)}
\KL\!\left(
q_{i-1|i}(\cdot\mid z_i,x)
\|p^\theta_{i-1|i}(\cdot\mid z_i)
\right).
\]
This proves the discrete-time ELBO stated on the previous slide.
\end{proofs}
\end{frame}

\begin{frame}{Why only $z_i$ remains outside the KL}
\scriptsize
For the $i$-th reverse term, start from the path expectation:
\[
\E_{q(z_{0:K}\mid x)}
\left[
\log
\frac{p^\theta_{i-1|i}(z_{i-1}\mid z_i)}
{q_{i-1|i}(z_{i-1}\mid z_i,x)}
\right].
\]
The integrand depends only on $(z_{i-1},z_i)$, so marginalize all other variables:
\[
\sum_{z_{i-1},z_i}q(z_{i-1},z_i\mid x)
\log
\frac{p^\theta_{i-1|i}(z_{i-1}\mid z_i)}
{q_{i-1|i}(z_{i-1}\mid z_i,x)}.
\]
Now split the pair marginal as
\[
q(z_{i-1},z_i\mid x)
=q_i(z_i\mid x)q_{i-1|i}(z_{i-1}\mid z_i,x).
\]
\end{frame}

\begin{frame}{Why only $z_i$ remains outside the KL (Cont)}
\scriptsize
\begin{proofs}[\proofname\ (Cont)]
Substitute the pair factorization:
\[
\sum_{z_i}q_i(z_i\mid x)
\sum_{z_{i-1}}q_{i-1|i}(z_{i-1}\mid z_i,x)
\log
\frac{p^\theta_{i-1|i}(z_{i-1}\mid z_i)}
{q_{i-1|i}(z_{i-1}\mid z_i,x)}.
\]
The inner sum is the negative KL:
\[
\sum_{z_{i-1}}q_{i-1|i}(z_{i-1}\mid z_i,x)
\log
\frac{p^\theta_{i-1|i}(z_{i-1}\mid z_i)}
{q_{i-1|i}(z_{i-1}\mid z_i,x)}
\]
\[
=-\KL\!\left(
q_{i-1|i}(\cdot\mid z_i,x)
\|p^\theta_{i-1|i}(\cdot\mid z_i)
\right).
\]
Therefore only the outer marginal $q_i(z_i\mid x)$ remains outside the KL.
\end{proofs}
\end{frame}

\begin{frame}{Small-time model reverse kernel}
\small
Let the model reverse CTMC have rate $\Rhat_t^\theta(y,z)$ from reverse source $z$ to reverse destination $y$. Write
\[
b_t(y,z):=\Rhat_t^\theta(y,z).
\]
For $s=t-\Delta$,
\[
p^\theta_{s|t}(y\mid z)
=
\begin{cases}
1+\Delta b_t(z,z)+o(\Delta),& y=z,\\
\Delta b_t(y,z)+o(\Delta),& y\neq z,
\end{cases}
\]
where
\[
b_t(z,z)=-\sum_{y\neq z}b_t(y,z).
\]

\begin{alertblock}{What the ELBO will compare}
For each current state $z$, the true reverse rates $a_t(y,z)$ are compared with the model reverse rates $b_t(y,z)$ over all possible destinations $y$.
\end{alertblock}
\end{frame}

\begin{frame}{Single-step KL expansion}
\small
Use the true and model small-time kernels:
\[
q_{s|t}(y\mid z,x)=
\begin{cases}
1+\Delta a_t(z,z)+o(\Delta),&y=z,\\
\Delta a_t(y,z)+o(\Delta),&y\neq z,
\end{cases}
\]
\[
p^\theta_{s|t}(y\mid z)=
\begin{cases}
1+\Delta b_t(z,z)+o(\Delta),&y=z,\\
\Delta b_t(y,z)+o(\Delta),&y\neq z.
\end{cases}
\]

\begin{block}{First-order term}
\[
\KL_\Delta(z)
=\Delta\left[
a_t(z,z)-b_t(z,z)
+\sum_{y\neq z}a_t(y,z)
\log\frac{a_t(y,z)}{b_t(y,z)}
\right]+o(\Delta).
\]
\end{block}

\begin{alertblock}{Key point}
Each reverse-posterior KL is $O(\Delta)$, so the sum over $1/\Delta$ time steps becomes an integral.
\end{alertblock}
\end{frame}

\begin{frame}{Single-step KL proof: diagonal term}
\scriptsize
\begin{proof}
Let
\[
\KL_\Delta(z)=
\KL(q_{s|t}(\cdot\mid z,x)\|p^\theta_{s|t}(\cdot\mid z)).
\]
The diagonal contribution is
\[
(1+\Delta a_t(z,z)+o(\Delta))
\log
\frac{1+\Delta a_t(z,z)+o(\Delta)}
{1+\Delta b_t(z,z)+o(\Delta)}.
\]
Using $\log(1+\Delta c+o(\Delta))=\Delta c+o(\Delta)$, this equals
\[
(1+\Delta a_t(z,z)+o(\Delta))
\bigl[\Delta(a_t(z,z)-b_t(z,z))+o(\Delta)\bigr]
\]
\[
=\Delta(a_t(z,z)-b_t(z,z))+o(\Delta).
\]
\end{proof}
\end{frame}

\begin{frame}{Single-step KL proof: off-diagonal terms}
\scriptsize
\begin{proofs}[\proofname\ (Cont)]
For $y\neq z$, the contribution is
\[
(\Delta a_t(y,z)+o(\Delta))
\log
\frac{\Delta a_t(y,z)+o(\Delta)}
{\Delta b_t(y,z)+o(\Delta)}.
\]
The factor $\Delta$ cancels inside the logarithm, so
\[
\log
\frac{\Delta a_t(y,z)+o(\Delta)}
{\Delta b_t(y,z)+o(\Delta)}
=\log\frac{a_t(y,z)}{b_t(y,z)}+o(1).
\]
Therefore the off-diagonal sum is
\[
\Delta\sum_{y\neq z}a_t(y,z)
\log\frac{a_t(y,z)}{b_t(y,z)}
+o(\Delta).
\]
Combining with the diagonal term proves the single-step expansion.
\end{proofs}
\end{frame}

\begin{frame}{From the KL expansion to the CT-ELBO}
\scriptsize
\begin{proofs}[\proofname\ (Cont)]
Let $t_i-t_{i-1}=\Delta$. The discrete ELBO contains
\[
\sum_{i=1}^K\E_{z_i\sim q_{t_i}}\KL_\Delta(z_i).
\]
Using the first-order expansion,
\[
\KL_\Delta(z_i)
=\Delta F_{t_i}(z_i)+o(\Delta),
\]
where
\[
F_t(z)=a_t(z,z)-b_t(z,z)
+\sum_{y\neq z}a_t(y,z)\log\frac{a_t(y,z)}{b_t(y,z)}.
\]
Thus
\[
\sum_{i=1}^K\E_{q_{t_i}}\KL_\Delta
\to
\int_0^1\E_{z\sim q_t}F_t(z)\,\dd t.
\]
Substituting this limit into the discrete ELBO gives the CT-ELBO.
\end{proofs}
\end{frame}

\begin{frame}{CT-ELBO in rate form}
\small
Taking the continuous-time limit gives
\[
\log p_\theta(x)\ge C_{\log}
-\int_0^1
\E_{z\sim q_t(\cdot\mid x)}
\left[
a_t(z,z)-b_t(z,z)
+\sum_{y\neq z}a_t(y,z)
\log\frac{a_t(y,z)}{b_t(y,z)}
\right]\dd t.
\]

\begin{block}{Trainable negative CT-ELBO}
After removing terms independent of $\theta$,
\[
\mathcal L_{\mathrm{CT}}(\theta)
=\int_0^1\E_{z\sim q_t}
\left[
\sum_{y\neq z}b_t(y,z)
-\sum_{y\neq z}a_t(y,z)\log b_t(y,z)
\right]\dd t+\const.
\]
\end{block}
\end{frame}

\begin{frame}{Standard forward/backward-rate form}
\scriptsize
Substitute
\[
a_t(y,z)=R_t(z,y)\frac{q_t(y\mid x)}{q_t(z\mid x)},
\qquad
b_t(y,z)=\Rhat_t^\theta(y,z).
\]
Then
\[
\log p_\theta(x)\ge C_{\log}+\int_0^1
\E_{z\sim q_t(\cdot\mid x)}
\biggl[
\Rhat_t^\theta(z,z)-R_t(z,z)
\]
\[
\qquad
+\sum_{y\neq z}R_t(z,y)\frac{q_t(y\mid x)}{q_t(z\mid x)}
\log
\frac{\Rhat_t^\theta(y,z)q_t(z\mid x)}
{R_t(z,y)q_t(y\mid x)}
\biggr]\dd t.
\]

\begin{alertblock}{Likelihood-bound view}
\begin{itemize}
  \item If SEDD chooses $\Rhat_t^\theta(y,z)=R_t(z,y)s_\theta(z,t)_y$, CT-ELBO recovers a diffusion-weighted score entropy objective.
  \item GIDD chooses $\Rhat_t^\theta(y,z)=R_t(z,y)q_t(y\mid x_\theta)/q_t(z\mid x_\theta)$.
\end{itemize}
\end{alertblock}
\end{frame}

% ======================================================================
\section{SEDD}
% ======================================================================

\begin{frame}{SEDD: modeling the reverse rate by ratios}
\small
SEDD learns the concrete score
\[
s_\theta(z,t)_y\approx \frac{p_t(y)}{p_t(z)},
\qquad y\neq z.
\]
The model reverse rate is
\[
\Rhat_t^\theta(y,z)=R_t(z,y)s_\theta(z,t)_y,
\qquad y\neq z.
\]

\begin{block}{Why this is natural}
The true reverse rate is
\[
a_t(y,z)=R_t(z,y)\frac{p_t(y)}{p_t(z)}.
\]
Thus SEDD separates the known forward edge $R_t(z,y)$ from the unknown probability ratio. This edge is the forward transition from the earlier candidate $y$ to the current token $z$.
\end{block}

\begin{alertblock}{Sequence interpretation}
$s_\theta(Z_t,t)_{i,y}$ estimates the ratio between the current sequence and the sequence obtained by replacing token $i$ by $y$.
\end{alertblock}
\end{frame}

\begin{frame}{SEDD forward processes}
\small
Take
\[
R_t=\dot\sigma(t)G_\pi,
\qquad
G_\pi=\Pi_\pi-I.
\]
For a clean token $x_0$,
\[
p_{t|0}(\cdot\mid x_0)
=\exp(\sigma(t)G_\pi)e_{x_0}.
\]

\begin{block}{Closed form}
Using the eigen-decomposition from the CTMC backbone,
\[
p_{t|0}(\cdot\mid x_0)
=e^{-\sigma(t)}e_{x_0}
+(1-e^{-\sigma(t)})\pi.
\]
\end{block}

\begin{alertblock}{Examples}
\begin{itemize}
  \item Uniform SEDD: choose $\pi$ uniform on $\cV$.
  \item Absorbing/MASK SEDD: choose $\pi=\delta_m$.
\end{itemize}
\end{alertblock}
\end{frame}

\begin{frame}{Score entropy}
\scriptsize
For a fully supported distribution $p$, weights $w_{xy}\ge0$, and positive scores $s_\theta(x)_y>0$, define
\[
L_{SE}(\theta;p,w)=
\E_{x\sim p}
\sum_{y\neq x}w_{xy}
\left[
s_\theta(x)_y-\frac{p(y)}{p(x)}\log s_\theta(x)_y
 +K\!\left(\frac{p(y)}{p(x)}\right)
\right],
\]
where
\[
K(r)=r(\log r-1).
\]

\begin{block}{Nonnegative pointwise gap}
Let $r=p(y)/p(x)$. Then
\[
s-r\log s+K(r)
=r\left(\frac{s}{r}-1-\log\frac{s}{r}\right)\ge0.
\]
\end{block}

\begin{alertblock}{Consistent minimizer}
The unique minimizer is
\[
s_\theta^*(x)_y=\frac{p(y)}{p(x)}.
\]
\end{alertblock}
\end{frame}

\begin{frame}{From implicit to denoising score entropy}
\small
The implicit form contains an unknown ratio:
\[
\E_{x\sim p}\sum_{y\neq x}w_{xy}\frac{p(y)}{p(x)}
\log s_\theta(x)_y.
\]
If $p$ is generated by corrupting clean data,
\[
p(x)=\sum_{x_0}p(x\mid x_0)p_0(x_0),
\]
then the denoising identity gives
\[
\E_{x\sim p}\sum_{y\neq x}\frac{p(y)}{p(x)}f(x,y)
=
\E_{x_0\sim p_0,\;x\sim p(\cdot\mid x_0)}
\sum_{y\neq x}\frac{p(y\mid x_0)}{p(x\mid x_0)}f(x,y).
\]

\begin{block}{Benefit}
The conditional ratio $p(y\mid x_0)/p(x\mid x_0)$ is computable from the forward corruption kernel, while $s_\theta(x)_y$ comes from one network forward pass.
\end{block}
\end{frame}

\begin{frame}{Standalone SEDD loss: denoising score entropy}
\scriptsize
The SEDD loss can be derived without any likelihood-bound argument.
Apply score entropy to the noisy marginal $p_t$ and use the denoising identity:
\[
L_{DSE}(\theta;t)
=
\E_{x_0\sim p_0,\;z\sim p_{t|0}(\cdot\mid x_0)}
\sum_{y\neq z}w_t(z,y)
\left[
s_\theta(z,t)_y
-\frac{p_{t|0}(y\mid x_0)}{p_{t|0}(z\mid x_0)}
\log s_\theta(z,t)_y
\right]+\const.
\]

\begin{block}{What has been used}
Only two ingredients are needed:
\begin{enumerate}
  \item score entropy is minimized at the probability ratio $p_t(y)/p_t(z)$;
  \item the denoising identity replaces the unknown marginal ratio by the computable conditional ratio.
\end{enumerate}
\end{block}
\end{frame}

\begin{frame}{Diffusion-weighted DSE}
\scriptsize
For a diffusion process, it is natural to weight only edges that can occur under the forward CTMC. With the column-vector convention, choose
\[
w_t(z,y)=R_t(z,y).
\]
Then the standalone denoising score-entropy objective becomes
\[
L_{DWDSE}(x_0)
=\int_0^T\E_{z\sim p_{t|0}}
\sum_{y\neq z}R_t(z,y)
\left[
s_\theta(z,t)_y
-\frac{p_{t|0}(y\mid x_0)}{p_{t|0}(z\mid x_0)}
\log s_\theta(z,t)_y
\right]\dd t+\const.
\]

\begin{alertblock}{Important point}
This is already a valid SEDD training objective from score-entropy ratio estimation. CT-ELBO is a separate route that explains when the same weighted objective is also the trainable part of a likelihood lower bound.
\end{alertblock}
\end{frame}

\begin{frame}{CT-ELBO connection for DWDSE}
\small
Condition on clean $x_0$. The true and model reverse rates are
\[
a_t(y,z)=R_t(z,y)
\frac{p_{t|0}(y\mid x_0)}{p_{t|0}(z\mid x_0)},
\qquad
b_t(y,z)=R_t(z,y)s_\theta(z,t)_y.
\]
Substituting into the trainable CT-NELBO gives, up to constants,
\[
L_{DWDSE}(x_0)
=\int_0^T\E_{z\sim p_{t|0}}
\sum_{y\neq z}R_t(z,y)
\left[
s_\theta(z,t)_y
-\frac{p_{t|0}(y\mid x_0)}{p_{t|0}(z\mid x_0)}
\log s_\theta(z,t)_y
\right]\dd t.
\]

\begin{alertblock}{How to read the connection}
Score entropy gives SEDD a ratio-estimation loss directly. CT-ELBO shows that choosing the diffusion weight $R_t(z,y)$ makes that loss coincide with the rate-level likelihood bound, up to constants.
\end{alertblock}
\end{frame}

\begin{frame}{SEDD sequence training}
\small
Sample clean sequences $X_0\sim p_{\mathrm{data}}$ and time $t\sim U(0,T)$. Then corrupt independently:
\[
X_t^{(i)}\sim p_{t|0}^{\mathrm{tok}}(\cdot\mid X_0^{(i)}).
\]
The network outputs positive scores
\[
s_\theta(X_t,t)\in\mathbb R_+^{L\times N}.
\]

\begin{block}{Single-sample sequence loss}
\[
\widehat L_{DWDSE}^{\mathrm{seq}}
=\dot\sigma(t)\sum_{i=1}^L\sum_{y\neq X_t^{(i)}}
G(X_t^{(i)},y)
\left[
s_\theta(X_t,t)_{i,y}
-\frac{p_{t|0}^{\mathrm{tok}}(y\mid X_0^{(i)})}
{p_{t|0}^{\mathrm{tok}}(X_t^{(i)}\mid X_0^{(i)})}
\log s_\theta(X_t,t)_{i,y}
\right].
\]
\end{block}
\end{frame}

\begin{frame}{SEDD sampling: Euler and tau-leaping}
\small
For a single token, reverse Euler from $t$ to $s=t-\Delta t$ uses
\[
\Prob_\theta(X_s=y\mid X_t=z)\approx
\begin{cases}
\Delta t\,R_t(z,y)s_\theta(z,t)_y,&y\neq z,\\
1-\Delta t\sum_{u\neq z}R_t(z,u)s_\theta(z,t)_u,&y=z.
\end{cases}
\]

\begin{block}{Tau-leaping for sequences}
For each position $i$,
\[
\pi_i(y)=
\begin{cases}
\Delta t\,R_t(z_t^{(i)},y)s_\theta(Z_t,t)_{i,y},&y\neq z_t^{(i)},\\
1-\Delta t\sum_{u\neq z_t^{(i)}}R_t(z_t^{(i)},u)s_\theta(Z_t,t)_{i,u},&y=z_t^{(i)}.
\end{cases}
\]
After clipping and normalization, sample positions independently.
\end{block}

\begin{alertblock}{Practical meaning}
Euler is rate-based and local in time. Tau-leaping trades exact one-jump CTMC semantics for many parallel token updates.
\end{alertblock}
\end{frame}

\begin{frame}{Discrete Tweedie update}
\small
If the forward process is $p_t=e^{\sigma(t)G}p_0$ and we know all ratios
\[
r_z(i)=\frac{p_t(i)}{p_t(z)},
\]
Bayes' rule and $p_0=e^{-\sigma(t)G}p_t$ give the one-token identity
\[
p_{0|t}(x_0\mid z)
=
\left[e^{-\sigma(t)G}r_z\right]_{x_0}
\left[e^{\sigma(t)G}\right]_{z,x_0}.
\]

\begin{block}{Tweedie tau-leaping}
For a step of size $\Delta\sigma_t=\sigma(t)-\sigma(t-\Delta t)$, set $r_i(y)=s_\theta(Z_t,t)_{i,y}$ and $r_i(z_t^{(i)})=1$. Then
\[
\tilde\pi_i(y)
=
\left[e^{-\Delta\sigma_tG}r_i\right]_y
\left[e^{\Delta\sigma_tG}\right]_{z_t^{(i)},y}.
\]
Normalize $\tilde\pi_i$ and sample each position.
\end{block}
\end{frame}

\begin{frame}{Why SEDD supports infilling}
\small
Let $\Omega$ be positions to generate and $\Omega^c$ be fixed prompt positions.
For two candidate fills $z,z'$,
\[
\frac{p_t(X_\Omega=z'\mid X_{\Omega^c}=y)}
{p_t(X_\Omega=z\mid X_{\Omega^c}=y)}
=
\frac{p_t(X_\Omega=z',X_{\Omega^c}=y)}
{p_t(X_\Omega=z,X_{\Omega^c}=y)}.
\]

\begin{block}{Consequence}
The unconditional score ratio is already the conditional score ratio for changes inside $\Omega$.
\end{block}

\begin{alertblock}{Inference rule}
Keep prompt positions fixed. Apply Euler or Tweedie tau-leaping only to the generated positions.
\end{alertblock}
\end{frame}

% ======================================================================
\section{GIDD}
% ======================================================================

\begin{frame}{GIDD: interpolating marginals}
\small
GIDD starts by specifying the marginal forward distribution:
\[
q_t(z\mid x)=\alpha_t\delta_{z,x}+\beta_t\pi_t(z),
\qquad
\beta_t=1-\alpha_t.
\]
Here
\[
\alpha_0=1,\qquad \alpha_1=0,\qquad \pi_t\in\DeltaV.
\]

\begin{block}{Interpretation}
At time $t$, a token is a linear interpolation between the clean one-hot token and a possibly time-varying mixing distribution.
\end{block}

\begin{alertblock}{Special case}
If $\pi_t=\mvec$ is the MASK point mass, GIDD reduces to masked diffusion.
\end{alertblock}
\end{frame}

\begin{frame}{Constructing forward transition probabilities}
\small
For $0\le s\le t\le1$, define
\[
\alpha_{t|s}:=\frac{\alpha_t}{\alpha_s},
\qquad
v_{t|s}(b):=\beta_t\pi_t(b)-\frac{\alpha_t}{\alpha_s}\beta_s\pi_s(b).
\]
Then set
\[
\Prob(X_t=b\mid X_s=a)=\alpha_{t|s}\delta_{b,a}+v_{t|s}(b).
\]

\begin{block}{Equivalent mixture form}
If $\beta_{t|s}=1-\alpha_{t|s}$ and $\pi_{t|s}=v_{t|s}/\beta_{t|s}$,
\[
\Prob(X_t=b\mid X_s=a)=\alpha_{t|s}\delta_{b,a}+\beta_{t|s}\pi_{t|s}(b).
\]
\end{block}

\begin{alertblock}{Marginal check}
Setting $s=0$ gives
\[
\Prob(X_t=z\mid X_0=x)=\alpha_t\delta_{z,x}+\beta_t\pi_t(z)=q_t(z\mid x).
\]
\end{alertblock}
\end{frame}

\begin{frame}{Proof: transition probabilities are valid}
\scriptsize
\begin{proof}
First check normalization over destination states. For any source state $a$,
\[
\sum_b\Prob(X_t=b\mid X_s=a)
=\sum_b\left(\frac{\alpha_t}{\alpha_s}\delta_{a,b}
+\beta_t\pi_t(b)-\frac{\alpha_t}{\alpha_s}\beta_s\pi_s(b)\right).
\]
Since $\sum_b\pi_t(b)=\sum_b\pi_s(b)=1$,
\[
\sum_b\Prob(X_t=b\mid X_s=a)
=\frac{\alpha_t}{\alpha_s}
+\beta_t-\frac{\alpha_t}{\alpha_s}\beta_s
=\frac{\alpha_t}{\alpha_s}(1-\beta_s)+\beta_t
=\alpha_t+\beta_t=1.
\]
Under the admissibility condition $v_{t|s}(b)\ge0$, all entries are nonnegative.
Hence these probabilities define a valid Markov transition law.
\end{proof}
\end{frame}

\begin{frame}{Proof: marginal consistency}
\scriptsize
\begin{proofs}[\proofname\ (Cont)]
At $s=0$, the boundary conditions give $\alpha_0=1$ and $\beta_0=0$. Therefore
\[
\alpha_{t|0}=\alpha_t,
\qquad
v_{t|0}(z)=\beta_t\pi_t(z).
\]
Thus
\[
\Prob(X_t=z\mid X_0=x)
=\alpha_t\delta_{z,x}+\beta_t\pi_t(z)
=q_t(z\mid x).
\]
This proves that the constructed CTMC has exactly the prescribed GIDD marginal.
\end{proofs}
\end{frame}

\begin{frame}{Proof: Chapman--Kolmogorov equation}
\scriptsize
\begin{proofs}[\proofname\ (Cont)]
Let $s\le t\le u$ and write
\[
\Prob(X_t=c\mid X_s=a)=\alpha_{t|s}\delta_{c,a}+v_{t|s}(c),
\quad
\Prob(X_u=b\mid X_t=c)=\alpha_{u|t}\delta_{b,c}+v_{u|t}(b).
\]
Then
\begin{align*}
\sum_c\Prob(X_u=b\mid X_t=c)\Prob(X_t=c\mid X_s=a)
&=\alpha_{t|s}\alpha_{u|t}\delta_{a,b}
+\alpha_{t|s}v_{u|t}(b)\\
&\quad+\alpha_{u|t}v_{t|s}(b)
+v_{u|t}(b)\sum_cv_{t|s}(c).
\end{align*}
Because $\sum_cv_{t|s}(c)=1-\alpha_{t|s}$, this becomes
\[
\alpha_{t|s}\alpha_{u|t}\delta_{a,b}
+v_{u|t}(b)+\alpha_{u|t}v_{t|s}(b).
\]
\end{proofs}
\end{frame}

\begin{frame}{Proof: Chapman--Kolmogorov equation (Cont)}
\scriptsize
\begin{proofs}[\proofname\ (Cont)]
The first coefficient telescopes:
\[
\alpha_{t|s}\alpha_{u|t}
=\frac{\alpha_t}{\alpha_s}\frac{\alpha_u}{\alpha_t}
=\frac{\alpha_u}{\alpha_s}
=\alpha_{u|s}.
\]
For the vector term,
\[
v_{u|t}+\alpha_{u|t}v_{t|s}
=\left(\beta_u\pi_u-\frac{\alpha_u}{\alpha_t}\beta_t\pi_t\right)
+\frac{\alpha_u}{\alpha_t}
\left(\beta_t\pi_t-\frac{\alpha_t}{\alpha_s}\beta_s\pi_s\right).
\]
The middle terms cancel:
\[
v_{u|t}+\alpha_{u|t}v_{t|s}
=\beta_u\pi_u-\frac{\alpha_u}{\alpha_s}\beta_s\pi_s
=v_{u|s}.
\]
Therefore
\[
\sum_c\Prob(X_u=b\mid X_t=c)\Prob(X_t=c\mid X_s=a)
=\Prob(X_u=b\mid X_s=a).
\]
\end{proofs}
\end{frame}

\begin{frame}{Grid proof and telescoping}
\scriptsize
For a single grid step, define
\[
\dot\alpha_t=\frac{\alpha_{t+\Delta}}{\alpha_t},
\qquad
\dot v_t
=\beta_{t+\Delta}\pi_{t+\Delta}
-\frac{\alpha_{t+\Delta}}{\alpha_t}\beta_t\pi_t.
\]
The Chapman--Kolmogorov calculation gives
\[
\alpha_{t+\Delta|s}=\dot\alpha_t\alpha_{t|s},
\qquad
v_{t+\Delta|s}=\dot v_t+\dot\alpha_t v_{t|s}.
\]
Iterating,
\[
\alpha_{t|s}
=\prod_i\frac{\alpha_{t_{i+1}}}{\alpha_{t_i}}
=\frac{\alpha_t}{\alpha_s},
\]
and
\[
v_{t|s}
=\sum_i\frac{\alpha_t}{\alpha_{t_{i+1}}}
\left(\beta_{t_{i+1}}\pi_{t_{i+1}}
-\frac{\alpha_{t_{i+1}}}{\alpha_{t_i}}\beta_{t_i}\pi_{t_i}\right)
=\beta_t\pi_t-\frac{\alpha_t}{\alpha_s}\beta_s\pi_s.
\]
\end{frame}

\begin{frame}{GIDD forward rate}
\small
Let $t\to t+\Delta$. Taylor expansion of the kernel gives
\[
\alpha_{t+\Delta|t}
=1+\frac{\alpha_t'}{\alpha_t}\Delta+o(\Delta),
\]
\[
v_{t+\Delta|t}(z)
=\left(\beta_t\pi_t'(z)-\frac{\alpha_t'}{\alpha_t}\pi_t(z)\right)\Delta+o(\Delta).
\]

\begin{block}{Forward rate}
\[
R_t(z,y)=\frac{\alpha_t'}{\alpha_t}\delta_{z,y}+h_t(z),
\qquad
h_t(z):=\beta_t\pi_t'(z)-\frac{\alpha_t'}{\alpha_t}\pi_t(z).
\]
\end{block}

\begin{alertblock}{Important structure}
For source $y$ and destination $z$ with $y\neq z$, the rate $R_t(z,y)=h_t(z)$ depends on the destination $z$, not on the source $y$.
\end{alertblock}
\end{frame}

\begin{frame}{GIDD backward parameterization}
\small
GIDD predicts a clean token distribution
\[
x_\theta(Z_t,t)\in\DeltaV.
\]
It then defines the model-induced marginal
\[
q_t(z\mid x_\theta)=\alpha_t x_\theta(z)+\beta_t\pi_t(z).
\]

\begin{block}{Posterior-form backward transition}
\[
p_\theta(z_s=y\mid z_t=z)
=\Prob(X_t=z\mid X_s=y)
\frac{q_s(y\mid x_\theta)}{q_t(z\mid x_\theta)}.
\]
\end{block}

\begin{alertblock}{Backward rate}
For $y\neq z$,
\[
\Rhat_t^\theta(y,z)
=R_t(z,y)
\frac{q_t(y\mid x_\theta)}{q_t(z\mid x_\theta)}.
\]
The diagonal is determined by the total rate leaving the current reverse state.
\end{alertblock}
\end{frame}

\begin{frame}{Proof: GIDD backward rate}
\scriptsize
\begin{proof}
Set $s=t-\Delta$. From the posterior-form backward transition,
\[
p_\theta(z_s=y\mid z_t=z)
=\Prob(X_t=z\mid X_s=y)\frac{q_s(y\mid x_\theta)}{q_t(z\mid x_\theta)}.
\]
For $y\neq z$,
\[
\Prob(X_t=z\mid X_s=y)=\Delta R_t(z,y)+o(\Delta),
\qquad
q_s(y\mid x_\theta)=q_t(y\mid x_\theta)+o(1).
\]
Therefore
\[
p_\theta(z_s=y\mid z_t=z)
=\Delta R_t(z,y)
\frac{q_t(y\mid x_\theta)}{q_t(z\mid x_\theta)}
+o(\Delta).
\]
Comparing with
\[
p_\theta(z_s=y\mid z_t=z)
=\Delta\Rhat_t^\theta(y,z)+o(\Delta)
\]
gives the off-diagonal backward rate.
\end{proof}
\end{frame}

\begin{frame}{Proof: diagonal backward rate}
\scriptsize
\begin{proofs}[\proofname\ (Cont)]
For $y=z$,
\[
p_\theta(z_s=z\mid z_t=z)
=1+\Delta\Rhat_t^\theta(z,z)+o(\Delta).
\]
Since the backward kernel must sum to one,
\[
0=\sum_y\Rhat_t^\theta(y,z)
=\Rhat_t^\theta(z,z)
+\sum_{y\neq z}R_t(z,y)
\frac{q_t(y\mid x_\theta)}{q_t(z\mid x_\theta)}.
\]
Hence
\[
\Rhat_t^\theta(z,z)
=-\sum_{y\neq z}R_t(z,y)
\frac{q_t(y\mid x_\theta)}{q_t(z\mid x_\theta)}.
\]
This is the diagonal part used in the CT-NELBO derivation.
\end{proofs}
\end{frame}

\begin{frame}{GIDD CT-NELBO proof: setup}
\scriptsize
\begin{proof}
Write
\[
q_z=q_t(z\mid x),\quad
\tilde q_z=q_t(z\mid x_\theta),\quad
h_z=h_t(z),\quad
w_z=\frac{h_z}{q_z}.
\]
For $y\neq z$, the GIDD forward rate and model reverse rate are
\[
R_t(z,y)=h_z,
\qquad
\Rhat_t^\theta(y,z)=h_z\frac{\tilde q_y}{\tilde q_z}.
\]
The CT-ELBO integrand for a fixed current state $z$ is
\[
I(z)=\Rhat_t^\theta(z,z)-R_t(z,z)
+\sum_{y\neq z}h_z\frac{q_y}{q_z}
\log
\frac{\Rhat_t^\theta(y,z)q_z}
{h_zq_y}.
\]
\end{proof}
\end{frame}

\begin{frame}{GIDD CT-NELBO proof: diagonal term}
\scriptsize
\begin{proofs}[\proofname\ (Cont)]
From the forward rate,
\[
R_t(z,z)=\frac{\alpha_t'}{\alpha_t}+h_z.
\]
The model diagonal rate is fixed by the total leaving rate:
\[
\Rhat_t^\theta(z,z)
=-\sum_{y\neq z}h_z\frac{\tilde q_y}{\tilde q_z}
=-h_z\frac{1-\tilde q_z}{\tilde q_z}
=h_z-\frac{h_z}{\tilde q_z}.
\]
Therefore
\[
\Rhat_t^\theta(z,z)-R_t(z,z)
=-\frac{h_z}{\tilde q_z}-\frac{\alpha_t'}{\alpha_t}.
\]
\end{proofs}
\end{frame}

\begin{frame}{GIDD CT-NELBO proof: logarithmic term}
\tiny
\begin{proofs}[\proofname\ (Cont)]
For $y\neq z$,
\[
\log
\frac{\Rhat_t^\theta(y,z)q_z}{h_zq_y}
=
\log\frac{\tilde q_y}{q_y}
+\log\frac{q_z}{\tilde q_z}.
\]
Hence
\begin{align*}
&\sum_{y\neq z}h_z\frac{q_y}{q_z}
\log
\frac{\Rhat_t^\theta(y,z)q_z}{h_zq_y}\\
&=w_z\sum_{y\neq z}q_y
\left(
\log\frac{\tilde q_y}{q_y}
+\log\frac{q_z}{\tilde q_z}
\right).
\end{align*}
Since
\[
\sum_{y\neq z}q_y\log\frac{\tilde q_y}{q_y}
=-\KL(q\|\tilde q)+q_z\log\frac{\tilde q_z}{q_z},
\]
the whole logarithmic term is
\[
w_z\left[-\KL(q\|\tilde q)+\log\frac{q_z}{\tilde q_z}\right].
\]
\end{proofs}
\end{frame}

\begin{frame}{GIDD CT-NELBO proof: combine terms}
\scriptsize
\begin{proofs}[\proofname\ (Cont)]
Combine the diagonal term and the logarithmic term:
\[
I(z)
=-\frac{\alpha_t'}{\alpha_t}
-\frac{h_z}{\tilde q_z}
+w_z\left[-\KL(q\|\tilde q)+\log\frac{q_z}{\tilde q_z}\right].
\]
Because $h_z=w_zq_z$,
\[
\frac{h_z}{\tilde q_z}=w_z\frac{q_z}{\tilde q_z}.
\]
Thus
\[
I(z)
=-\frac{\alpha_t'}{\alpha_t}
-w_z\left[
\KL(q\|\tilde q)
+\frac{q_z}{\tilde q_z}
-\log\frac{q_z}{\tilde q_z}
\right].
\]
\end{proofs}
\end{frame}

\begin{frame}{GIDD CT-NELBO proof: weight identity}
\scriptsize
\begin{proofs}[\proofname\ (Cont)]
The weight has expectation
\[
\E_{z\sim q}[w_z]
=\sum_zq_z\frac{h_z}{q_z}
=\sum_zh_z.
\]
Using
\[
h_z=\beta_t\pi_t'(z)-\frac{\alpha_t'}{\alpha_t}\pi_t(z),
\]
we get
\[
\sum_zh_z
=\beta_t\frac{\dd}{\dd t}\sum_z\pi_t(z)
-\frac{\alpha_t'}{\alpha_t}\sum_z\pi_t(z)
=-\frac{\alpha_t'}{\alpha_t}.
\]
Taking $-\E_{z\sim q}[I(z)]$ and substituting this identity gives
\[
\E_{z\sim q}w_z
\left[
\KL(q\|\tilde q)
+\frac{q_z}{\tilde q_z}
-\log\frac{q_z}{\tilde q_z}
-1
\right].
\]
\end{proofs}
\end{frame}

\begin{frame}{GIDD CT-NELBO proof: final loss}
\scriptsize
\begin{proofs}[\proofname\ (Cont)]
By definition,
\[
\DIS(a\|b)=\frac{a}{b}-\log\frac{a}{b}-1.
\]
Therefore the negative CT-ELBO integrand becomes
\[
w_t(z,x)\left[
\KL(q_t(\cdot\mid x)\|q_t(\cdot\mid x_\theta))
+\DIS(q_t(z\mid x)\|q_t(z\mid x_\theta))
\right],
\]
where
\[
w_t(z,x)=
\frac{1}{q_t(z\mid x)}
\left(\beta_t\pi_t'(z)-\frac{\alpha_t'}{\alpha_t}\pi_t(z)\right).
\]
Integrating over $t$ and $z\sim q_t(\cdot\mid x)$ gives
\[
-\log p_\theta(x)
\le
\E_{t,\;z\sim q_t(\cdot\mid x)}
\left[
w_t(z,x)
\left(
\KL(q_t(\cdot\mid x)\|q_t(\cdot\mid x_\theta))
+\DIS(q_t(z\mid x)\|q_t(z\mid x_\theta))
\right)
\right]+C.
\]
\end{proofs}
\end{frame}

\begin{frame}{Mask-only special case}
\small
Let $\pi_t=\mvec$ be the MASK point mass. Then $\pi_t'=0$ and
\[
q_t(m\mid x)=1-\alpha_t.
\]
The weight becomes
\[
w_t(z,x)=
-\frac{\alpha_t'}{\alpha_t(1-\alpha_t)}\delta_{z,m}.
\]

\begin{block}{The loss reduces to masked reconstruction}
When $z=m$,
\[
\KL(q_t(\cdot\mid x)\|q_t(\cdot\mid x_\theta))
+\DIS(q_t(m\mid x)\|q_t(m\mid x_\theta))
=-\alpha_t\log x_\theta(x).
\]
Therefore
\[
-\log p_\theta(x)
\le
\E_{t,z_t}
\left[
\frac{\alpha_t'}{1-\alpha_t}\delta_{z_t,m}
\log x_\theta(x\mid Z_t,t)
\right]+C.
\]
\end{block}
\end{frame}

\begin{frame}{Mask plus uniform schedule}
\small
GIDD can keep the final prior as all MASK while injecting uniform noise in the middle:
\[
\alpha_t=\frac{1-t}{C_t},
\qquad
\beta_t\pi_t
=\frac{t}{C_t}\mvec+\frac{c_t}{C_t}\uvec,
\qquad
C_t=1+c_t.
\]
Here
\[
\uvec=\frac{1}{N-1}(\one-\mvec),
\qquad
c_t=B\,t^{\gamma/2}(1-t)^{\gamma/2}.
\]

\begin{block}{Choosing the maximum uniform ratio}
If the desired maximum uniform mass is $p_u$, choose
\[
B=2^\gamma\frac{p_u}{1-p_u}.
\]
Then $c_t/(1+c_t)$ is maximized at $t=1/2$ with value $p_u$.
\end{block}
\end{frame}

\begin{frame}{Closed-form weights for mask plus uniform}
\scriptsize
For the mask plus uniform schedule,
\[
c_t'=\frac{\gamma}{2}\frac{1-2t}{t(1-t)}c_t,
\qquad
\frac{\alpha_t'}{\alpha_t}
=-\frac{1}{1-t}-\frac{c_t'}{1+c_t}.
\]
The numerator of the ELBO weight simplifies to
\[
\beta_t\pi_t'-\frac{\alpha_t'}{\alpha_t}\pi_t
=\frac{\mvec+(c_t+(1-t)c_t')\uvec}{C_t(1-t)}.
\]
Hence
\[
w_t(z,x)
=e_z^\top
\left[
\frac{\mvec+(c_t+(1-t)c_t')\uvec}
{(1-t)((1-t)\xvec+t\mvec+c_t\uvec)}
\right].
\]

\begin{alertblock}{Training note}
The exact ELBO weights can be large near the endpoints. Practical GIDD variants often clip or dynamically reweight them to improve gradient stability.
\end{alertblock}
\end{frame}

\begin{frame}{GIDD sequence training}
\small
For each clean sequence $X=(x^{(1)},\ldots,x^{(L)})$:
\[
Z_t^{(i)}\sim
\Cat\!\left(\alpha_t e_{x^{(i)}}+\beta_t\pi_t\right).
\]
The network predicts
\[
x_\theta^{(i)}(Z_t,t)\in\DeltaV.
\]
Define
\[
q_t^{(i)}=\alpha_t e_{x^{(i)}}+\beta_t\pi_t,
\qquad
q_{\theta,t}^{(i)}=\alpha_t x_\theta^{(i)}+\beta_t\pi_t.
\]

\begin{block}{Sequence loss}
\[
\widehat L_{\mathrm{GIDD}}^{\mathrm{seq}}
=\sum_{i=1}^L
w_t(z_t^{(i)},x^{(i)})
\left[
\KL(q_t^{(i)}\|q_{\theta,t}^{(i)})
+\DIS(q_t^{(i)}(z_t^{(i)})\|q_{\theta,t}^{(i)}(z_t^{(i)}))
\right].
\]
\end{block}
\end{frame}

\begin{frame}{GIDD sampling}
\small
Choose a grid
\[
0\approx t_0<t_1<\cdots<t_K\approx1.
\]
Since the practical schedule has $\pi_1=\mvec$, initialize from the all-MASK sequence.

\begin{block}{Ancestral step}
For $s<t$,
\[
p_\theta(z_s=y\mid z_t=z)
=\Prob(X_t=z\mid X_s=y)
\frac{q_s(y\mid x_\theta(Z_t,t))}
{q_t(z\mid x_\theta(Z_t,t))}.
\]
For the forward transition probability,
\[
\Prob(X_t=z\mid X_s=y)
=\frac{\alpha_t}{\alpha_s}\delta_{z,y}
+\beta_t\pi_t(z)
-\frac{\alpha_t}{\alpha_s}\beta_s\pi_s(z).
\]
\end{block}

All positions can be updated in parallel after clipping numerical negatives and normalizing.
\end{frame}

\begin{frame}{Self-correction in GIDD}
\scriptsize
The uniform component means the model sees non-MASK but incorrect tokens during training. This enables post-denoising correction.

\begin{block}{One-token correction loop}
Given a nearly clean sequence $Z$, compute
\[
x_\theta^{(1:L)}=\softmax(f_\theta(Z,t)/\tau)
\]
at small $t$. Sample candidates $z'^{(i)}\sim\Cat(x_\theta^{(i)})$ and define
\[
S=\{i:z'^{(i)}\neq z^{(i)}\}.
\]
If $S$ is nonempty, update only
\[
j=\argmax_{i\in S}x_\theta^{(i)}(z'^{(i)}),
\qquad
z^{(j)}\leftarrow z'^{(j)}.
\]
\end{block}

\begin{alertblock}{Reason for one-at-a-time updates}
Changing one token reduces the risk of introducing multiple incompatible edits at once.
\end{alertblock}
\end{frame}

% ======================================================================
\section{Comparison}
% ======================================================================

\begin{frame}{SEDD and GIDD side by side}
\scriptsize
\begin{tabular}{@{}L{0.17\textwidth}L{0.36\textwidth}L{0.36\textwidth}@{}}
\toprule
Aspect & SEDD & GIDD \\
\midrule
Forward process &
General CTMC, often $R_t=\dot\sigma(t)G_\pi$ with uniform or MASK prior $\pi$. &
Specify $q_t=\alpha_tx+\beta_t\pi_t$, then construct a compatible CTMC. \\
\addlinespace
Network output &
Concrete score $s_\theta(z,t)_y\approx p_t(y)/p_t(z)$. &
Clean token distribution $x_\theta(Z_t,t)$. \\
\addlinespace
Reverse rate &
$\Rhat_t^\theta(y,z)=R_t(z,y)s_\theta(z,t)_y$. &
$\Rhat_t^\theta(y,z)=R_t(z,y)q_t(y\mid x_\theta)/q_t(z\mid x_\theta)$. \\
\addlinespace
Training &
Score entropy gives ratio estimation directly; DWDSE uses the reverse-edge forward rate. &
Closed-form CT-NELBO with KL plus pointwise IS divergence. \\
\addlinespace
Sampling &
Euler, tau-leaping, and Tweedie tau-leaping. &
Posterior-form ancestral sampling and optional self-correction. \\
\bottomrule
\end{tabular}
\end{frame}

\begin{frame}{The conceptual difference}
\small
\begin{columns}[T,totalwidth=\textwidth]
\column{0.48\textwidth}
\begin{block}{SEDD}
SEDD asks the model to learn the ratio that the reverse CTMC needs:
\[
\frac{p_t(y)}{p_t(z)}.
\]
This directly matches the true reverse rate identity and naturally supports conditional sampling through ratio cancellation.
\end{block}

\column{0.48\textwidth}
\begin{alertblock}{GIDD}
GIDD asks the model to predict the clean token distribution:
\[
x_\theta(Z_t,t).
\]
The reverse kernel is then obtained by plugging this prediction into the exact posterior formula for the chosen interpolating marginal.
\end{alertblock}
\end{columns}

\vspace{0.25cm}
Both use the same reverse-rate identity. SEDD has a standalone score-entropy derivation; CT-ELBO gives its likelihood-bound interpretation. GIDD is derived by plugging its posterior-form rate into CT-ELBO.
\end{frame}

\begin{frame}{Core formulas to remember}
\scriptsize
\begin{enumerate}
	  \item True reverse rate:
	  \[
	  a_t(y,z)=R_t(z,y)\frac{q_t(y\mid x)}{q_t(z\mid x)},\qquad y\neq z.
	  \]

  \item Trainable CT-NELBO:
  \[
  \int\E_{z\sim q_t}
  \left[\sum_{y\neq z}b_t(y,z)-\sum_{y\neq z}a_t(y,z)\log b_t(y,z)\right]\dd t+\const.
  \]

	  \item SEDD:
	  \[
	  b_t(y,z)=R_t(z,y)s_\theta(z,t)_y.
	  \]

  \item GIDD:
  \[
	  q_t(z\mid x)=\alpha_t\delta_{z,x}+\beta_t\pi_t(z),
	  \qquad
	  b_t(y,z)=R_t(z,y)\frac{q_t(y\mid x_\theta)}{q_t(z\mid x_\theta)}.
	  \]
\end{enumerate}
\end{frame}

\begin{frame}{Summary}
\scriptsize
\begin{block}{Shared foundation}
Discrete diffusion can be treated as a CTMC. The single most important identity is the reverse jump rate:
\[
a_t(y,z)=R_t(z,y)\frac{q_t(y\mid x)}{q_t(z\mid x)}.
\]
\end{block}

\begin{alertblock}{Model design choices}
\begin{itemize}
  \item SEDD learns the ratio in this identity and trains it with score entropy; CT-ELBO is an additional likelihood interpretation of the diffusion-weighted version.
  \item GIDD designs an interpolating marginal and predicts clean distributions, yielding a closed-form weighted CT-NELBO.
  \item Sequence implementations rely on token-wise corruption but context-aware denoisers.
\end{itemize}
\end{alertblock}

\begin{block}{Practical takeaway}
Once the forward CTMC and the reverse-rate parameterization are clear, the training losses and samplers are consequences rather than separate tricks.
\end{block}
\end{frame}

\begin{frame}[allowframebreaks]{References}
\small
\begin{thebibliography}{9}
\bibitem{lou2024sedd}
Aaron Lou, Chenlin Meng, and Stefano Ermon.
\newblock Discrete Diffusion Modeling by Estimating the Ratios of the Data Distribution.
\newblock ICML 2024.

\bibitem{vonrutte2025gidd}
Dimitri von R\"utte, Janis Fluri, Yuhui Ding, Antonio Orvieto, Bernhard Sch\"olkopf, and Thomas Hofmann.
\newblock Generalized Interpolating Discrete Diffusion.
\newblock ICML 2025.

\bibitem{campbell2022ct}
Andrew Campbell, Joe Benton, Valentin De Bortoli, Tom Rainforth, George Deligiannidis, and Arnaud Doucet.
\newblock A Continuous Time Framework for Discrete Denoising Models.
\newblock 2022.
\end{thebibliography}
\end{frame}

\end{document}
